\subsection{National Diet Library OCR stack}
\label{sec:appendix-digitization}

Scanned directory pages are processed with the National Diet Library (NDL) open-source document OCR (version 2.x), which chains several neural and classical steps before line recognition. The pipeline, controlled through a configuration file and a command-line driver, typically runs in four stages: (i) \emph{page separation}, which detects and splits two-up spreads when present; (ii) \emph{deskew}, which corrects small rotation; (iii) \emph{layout extraction}, which detects text regions on the page with a convolutional detector; and (iv) \emph{line OCR}, which transcribes each line with a Japanese-character recognizer trained on historical print. The system can emit reading-order information, ruby (furigana) handling, and optional dumps of intermediate crops. For archival batch jobs we run inference on GPU nodes and request XML output so that each line retains its text content together with coordinates (and page or image identifiers) for downstream parsing.

The implementation bundles pretrained weights for separation, layout, and recognition; line recognition uses a character inventory suited to kanji--kana mixtures found in prewar administrative volumes. Batch runs store images alongside machine-readable XML in a dedicated output folder. Image quality and script style vary by era: interwar calligraphy and wartime paper can reduce accuracy relative to postwar editions, which motivates explicit validation and noise flags in the compilation stage described below.

\subsection{From OCR XML to row-level personnel records}

The extraction stage treats NDL XML as the interchange format between OCR and statistical analysis. For each publication year and government level (Tokyo city, Tokyo prefecture, or Tokyo metropolis), routines ingest all XML files for that wave. For each page they parse the XML tree, extract every text line with its bounding box, and cluster lines into vertical columns so that reading order follows traditional right-to-left columns and top-to-bottom order within a column---matching how names and titles appear in the printed directories.

Before splitting names, the code strips recurring metadata embedded in the same string as personal names: monthly salary expressed in kanji numerals, court rank strings, and grade or appointment markers. Remaining text is tokenized with SudachiPy (mode tuned for conservative splits). Tokens tagged as person names are retained; other material is kept for position matching and diagnostics. Job titles are matched against an external position crosswalk that maps variant strings to standardized labels used in analysis; unmatched titles are labeled as unknown or passed through as raw text depending on the rule set.

\subsection{Compilation: offices, positions, and noise flags}

A second stage converts the raw line-level table into an analysis-ready file for that year. It classifies lines as office headers using suffix patterns (e.g., bureau, division, section), entries in the crosswalk, and guards against confusing a position title (e.g., section chief) with an office name. Detected headers are forward-filled so that each employee row carries a nested administrative hierarchy (bureau through subsection, with normalized orthography mapping from historical to modern characters where applicable). The routine flags table-of-contents or index pages, library stamps, and other non-employee lines, and sets a boolean indicating whether a row plausibly represents a person's name (length, character set, particle checks). Heuristic gender labels are assigned from historical Japanese naming patterns (including conservative and extended rule sets). Drafting and wartime status markers are applied where encoded in the source text.

\subsection{Master panel and identifiers}

After each year-level file is produced, we concatenate all editions into a long panel. Each distinct office string receives a deterministic integer identifier for use in fixed-effects specifications. Cross-year person identifiers are assigned on name rows by grouping identical normalized names and, where implemented, applying similarity thresholds within offices to merge likely duplicate transcriptions; short names may be restricted to exact matching to limit false merges. Page number, image filename, and line coordinates from OCR are retained for auditing and for linking back to scans.

Implementation details, software dependencies, and data documentation are available in the supplementary materials; large model weights and proprietary crosswalks may reside on secured storage.
